\section{A ternary-tree source theorem closes every four-or-more-factor crossing}
This is an unbounded analytic refinement of the constant-five entrance
reduction. It concerns late positive-derivative LC. General C, A2 and
original unamplified one-descent coefficient unimodality remain OPEN.
There is NO new finite-host entrance premise in this proof.

\paragraph{Convex-source subtraction after geometric tilt.}
Let a nonnegative LC background triple be $(X,x,Z)$, with
$x\le RX$, $R>0$. Suppose a nonnegative source $(e_-,e,e_+)$
satisfies $e\ge Re_-$ and
$e_+-2Re+R^2e_-\ge0$. Then $(X-e_-,x-e,Z-e_+)^+$ is LC.
If a neighbour is zero this is automatic. Otherwise the background
triple is positive, by its LC, so ratios below are legitimate.
Ordinary LC and $x\le RX$ imply $Z\le Rx$. Put
$C=R^2xX-2RXZ+xZ$. The exact identity
\[
 xC=X(Rx-Z)^2+Z(x^2-XZ)\ge0
\]
and the next identity prove convexity of the source density:
\[
 e_-xZ+e_+Xx-2eXZ
  =Xx(e_+-2Re+R^2e_-)+2X(Rx-Z)(e-Re_-)+e_-C\ge0.
\]
Thus $a=e_-/X,b=e/x,c=e_+/Z$ satisfy $a+c\ge2b$.
Their positive complementary triple is LC: if neighbours survive,
$b<1$ and its exact minor is
$(b-a)^2+(1-a)(a+c-2b)\ge0$.
Pointwise multiplication by the LC background proves the theorem.
The complete clipped scalar implication, not only its polynomial
identity, is kernel checked in \texttt{ConvexTiltLagrangeSource.lean}.

\paragraph{A uniformly positive ternary diagonal.}
Let $p=\phi$, $q=p-1$, and
$U^{(r)}_j=[z^j](1+z)^u/(1-z)^r$. Write
\[
 Q^{(r)}_{u,j}=U^{(r)}_{j+1}-2pU^{(r)}_j+p^2U^{(r)}_{j-1},
 \qquad Q^{(r)}_{u,-1}=1.
\]
For ALL $j\ge1$, $u\ge3j$, we prove
\[
                  Q^{(4)}_{u,j}\ge26-13p>4.          \tag{1}
\]
The standard external input is the formal Lagrange formula in
Ira M. Gessel, \emph{Lagrange Inversion}, Theorem 2.1.1,
equation (2.1.4),
\texttt{https://arxiv.org/html/1609.05988}; JCTA144(2016),212--249,
DOI \texttt{10.1016/j.jcta.2016.06.018}. No priority is claimed
for that formula. If $w=t(1+w)^3$, it states
\[
 [t^k]\frac{f(w)}{1-t\Phi'(w)}=[z^k]f(z)\Phi(z)^k,
                         \qquad\Phi(z)=(1+z)^3.
\]
For fixed integer $e\ge0$, substitute
$f(z)=(1-pz)^2(1+z)^{e-3}/(1-z)^4$.
Its diagonal series becomes
\[
 G(w)(1+w)^e,\qquad
 G(z)=\frac{(1-pz)^2}{(1-z)^2(1-z^2)^2(1-2z)}.
\]
At $k=j+1$, the original coefficient is exactly
$Q^{(4)}_{3j+e,j}$. This equality is valid in formal series
over $\mathbb Q(p)$; the negative exponent $e-3$ is permitted.

All coefficients of $G$ are positive. Indeed, for
$g(z)=(1-pz)^2/((1-z)^2(1-2z))$, one has
$g_0=1,g_1=4-2p$, and, for $n\ge2$,
$g_n=q^4 2^n-q^2n+3p-4$.
Here $g_2=12-7p>0$ and
$g_{n+1}-g_n=q^4 2^n-q^2\ge18-11p>0$, using $p\le13/8$.
The complete all-rank lower bound $g_n\ge5/8$ is kernel checked.
Multiplication by $(1-z^2)^{-2}$ preserves coefficient positivity;
specifically $G_1=4-2p$ and $G_2=14-7p$.
The unique ternary-tree series $w$ has positive integer coefficients,
with $[t^k]w\ge3$ and $[t^k]w^2\ge1$ for every $k\ge2$.
These follow directly from $w=t+3tw+3tw^2+tw^3$.
Since $(1+w)^e$ has nonnegative coefficients,
$[t^k]G(w)(1+w)^e\ge3G_1+G_2=26-13p$.
This proves (1) for EVERY $u=3j+e\ge3j$.

For EVERY $r\ge4$, multiplication by $(1-z)^{-(r-4)}$ now gives
\[
 Q^{(r)}_{u,j}\ge(r-3)(26-13p)>r
                  \qquad(j\ge2,\ u\ge3j).          \tag{2}
\]
Retain the first two convolution terms and use (1) at $j$ and $j-1$.
All remaining terms are nonnegative: they have source indices
$-1..j$, and $Q^{(4)}_{u,0}=u+4-2p>0$ as well.
The strict scalar payment of the possible $r$ unit cuts is kernel
checked for ALL real $r\ge4$.

\paragraph{The true CI source, including its first capacity cuts.}
For a deletion-minimal crossing write $h=u+\delta$, $\delta\ge3$,
$r\ge4$, and $d_i=t_i-1\ge\delta$.
Let $L=\prod[t_i]$, $E=(1+z)^uL$.
Through source rank $\delta$, $E$ is EXACTLY $U^{(r)}$:
every capacity bound is still inactive. At rank $\delta+1$ its
only losses are the individual unit cuts at capacities $\delta+1$,
in number at most $r$. Because $(r-1)\delta\le u$, every
$2\le j\le\delta$ satisfies $u\ge3j$. Equation (2) therefore
proves
\[
 E_{j+1}-2\phi E_j+\phi^2E_{j-1}>0
                          \quad(2\le j\le\delta).  \tag{3}
\]
Also $jE_j\ge(u-j+1)E_{j-1}$, so $E_j\ge2E_{j-1}$
on this range. These are precisely the convex-source conditions
against any already proved subgolden LC background.

\paragraph{The natural quotient entrance at every count.}
Counts $u\le13$ are already covered by the registered all-small-binary
crossing theorem. Assume $u\ge14$. Remove the LARGEST $[t]$.
Minimality and $r\ge4$ give
\[
 h_0=u+\delta-t+1\le u,\quad 3\le r_0=r-1\le u/3,
                           \quad t\le u-5.          \tag{4}
\]
Its old quotient is a registered positive global-LC degree-boundary
or surplus seed with later-primitive LC. No exceptional $u=3$
quotient occurs.
The ULC factor from the preceding unscaled theorem improves to
an UNSCALED UNIT interior kernel already at $u\ge12$:
\[
 K(s,k)\ge T_k\quad(2\le s\le u-3,\ 1\le k\le s).
\]
Its scalar margin, after $q\le2/3$, is
$3u^3-26u^2-26u+4$; at $u=12+x$ this is
$3x^3+82x^2+646x+1132>0$.
For every interval-supported LC numerator $N_0=1$,
$N_1\le u/3+1$, with NO degree bound, Cauchy--Binet gives
\[
 M_{NQ_u}(s)\ge B_{s+1}+2qB_s+q^2B_{s-1}
                                  \quad(2\le s\le u-3).
\]
Row zero is paid by $\binom{u-1}{2}\ge4u/3+5/3$.

The stronger norm for capacities at least four is
$\rho=q+q^6<27/40$, so for every $r\ge4$,
$\phi^2\rho^r\le(21/8)(27/40)^4<3/5$.
Consequently exactly the natural-head proof of the >=6-factor theorem
applies, with this $3/5$ atom and the unit budget:
the quotient tail from rank $\delta-1$ is positive global LC,
and its positive primitive after deleting rank zero is global LC.
All ordinary and primitive head centres are within (4)'s interior.
The only changed junction estimate is the unperturbed $3u/2$ gap.
It holds uniformly on $3\le r_0\le u/3$, $u\ge12$, by concavity;
its endpoint margins are
$(5u^2-49u+32)/2$ and $(2u^3-17u^2-12u+9)/9$.
The atom error is at most $4u/5+2/5$, leaving a strict margin.
All these unbounded scalar and interval-gap implications are checked
by the new Lean file. No finite entrance cases are substituted here.

\paragraph{All finite paddings and the constant-three conclusion.}
The geometric inverse of the positive LC natural quotient gives the
ACTUAL infinite-initial-ones signed derivative tail from physical
$W$, including its indispensable previous unit. Its ratios are
strictly less than $\phi$ and its pre-mode triples are LC.
For lag $a=\alpha-1$ the full finite source remains $E_{s-a}$,
including the spectral atom. Discarding original primitive indices
below $\delta$ is legitimate in EVERY output triple from centre
$a+\delta+1$, by the registered kernel-prefix support theorem.
The signed-primitive LC/constant-cut proof pays every later triple.
Before that range, source indices are at most $\delta$.
Ranks zero and one use their registered exact startup payments;
equation (3) and the convex-source theorem pay every rank $2..\delta$.
The old same-word derivative-density order and inverse-mode bound
put the finite first tail mode inside this proved range.

Thus EVERY deletion-minimal crossing with AT LEAST FOUR CI factors
has actual all-padding late positive-derivative LC. Together with
the already proved single-factor, degree-$u+1$ and degree-$u+2$
seeds and arbitrary CI insertion, the ENTIRE universal late-D-LC
entrance now needs only TWO- or THREE-factor minimal crossings:
\[
 r\in\{2,3\},\quad 3\le d_i\le u,\quad
 u+3\le h\le2u,\quad h-\min d_i\le u,
\]
and they may have no subset of degree $u+1$ or $u+2$.
Their binary count and capacities are still unbounded. No assertion
is made that these last entrances have been verified in full.
General C/A2/original UI remain OPEN; actual late-C threshold
$u\ge h-2$ is unchanged, and no original numerical bound is upgraded.

The formal Lagrange input and series coefficient/composition assembly
are written mathematics, supported by primary literature and independent
exact identities. Lean checks the new unbounded scalar, clipped-source,
all-rank kernel and gap implications. No end-to-end arbitrary-composition
or complete formal-power-series Lean theorem is claimed.
